%! Simplifying Radicals; Adding \& Subtracting — Unit 5, Lesson 5.2 — Exit Ticket
\documentclass[10pt]{article}
\usepackage{algebra2-article}
\usepackage{algebra2-boxes}
\newcommand{\TallMath}[1]{$\displaystyle #1\rule[-1.4em]{0pt}{3.2em}$}

\begin{document}

\pageheader{Unit 5, Lesson 5.2}{Exit Ticket}
\namedateperiod

\vspace{0.10in}

No notes, no calculator. Assume all variables represent positive real numbers.

\begin{enumerate}[leftmargin=1.8em, itemsep=10pt]
  \item \textbf{Simplify completely.}

        \vspace{4pt}
        (a) $\sqrt{75} = \blank{2.2cm}$ \hspace{0.5in}
        (b) $\sqrt[3]{54} = \blank{2.2cm}$

        \vspace{4pt}
        (c) $\sqrt{18x^{5}} = \blank{2.6cm}$ \hspace{0.4in}
        (d) $2\sqrt{12}+\sqrt{27} = \blank{2.6cm}$

  \item \textbf{Explain.} A classmate writes $\sqrt{9}+\sqrt{16}=\sqrt{25}=5$. Show with numbers
        that this is wrong, then state the rule that \emph{does} hold and say what operation it
        requires.
        \writelines{3}

  \item \textbf{SOL-style multiple choice.} Which expression is equivalent to
        $\sqrt{50}+\sqrt{32}$?
        \begin{enumerate}[label=\Alph*., leftmargin=1.9em, itemsep=1pt]
          \item $\sqrt{82}$
          \item $9\sqrt{2}$
          \item $9\sqrt{4}$
          \item $20\sqrt{2}$
        \end{enumerate}
\end{enumerate}

\end{document}
